\IfFileExists{stacks-project.cls}{%
\documentclass{stacks-project}
}{%
\documentclass{amsart}
}

% For dealing with references we use the comment environment
\usepackage{verbatim}
\newenvironment{reference}{\comment}{\endcomment}
%\newenvironment{reference}{}{}
\newenvironment{slogan}{\comment}{\endcomment}
\newenvironment{history}{\comment}{\endcomment}

% For commutative diagrams we use Xy-pic
\usepackage[all]{xy}

% We use 2cell for 2-commutative diagrams.
\xyoption{2cell}
\UseAllTwocells

% We use multicol for the list of chapters between chapters
\usepackage{multicol}

% This is generally recommended for better output
\usepackage{lmodern}
\usepackage[T1]{fontenc}

% For cross-file-references
\usepackage{xr-hyper}

% Package for hypertext links:
\usepackage{hyperref}

% For any local file, say "hello.tex" you want to link to please
% use \externaldocument[hello-]{hello}
\externaldocument[introduction-]{introduction}
\externaldocument[conventions-]{conventions}
\externaldocument[sets-]{sets}
\externaldocument[categories-]{categories}
\externaldocument[topology-]{topology}
\externaldocument[sheaves-]{sheaves}
\externaldocument[sites-]{sites}
\externaldocument[stacks-]{stacks}
\externaldocument[fields-]{fields}
\externaldocument[algebra-]{algebra}
\externaldocument[brauer-]{brauer}
\externaldocument[homology-]{homology}
\externaldocument[derived-]{derived}
\externaldocument[simplicial-]{simplicial}
\externaldocument[more-algebra-]{more-algebra}
\externaldocument[smoothing-]{smoothing}
\externaldocument[modules-]{modules}
\externaldocument[sites-modules-]{sites-modules}
\externaldocument[injectives-]{injectives}
\externaldocument[cohomology-]{cohomology}
\externaldocument[sites-cohomology-]{sites-cohomology}
\externaldocument[dga-]{dga}
\externaldocument[dpa-]{dpa}
\externaldocument[sdga-]{sdga}
\externaldocument[hypercovering-]{hypercovering}
\externaldocument[schemes-]{schemes}
\externaldocument[constructions-]{constructions}
\externaldocument[properties-]{properties}
\externaldocument[morphisms-]{morphisms}
\externaldocument[coherent-]{coherent}
\externaldocument[divisors-]{divisors}
\externaldocument[limits-]{limits}
\externaldocument[varieties-]{varieties}
\externaldocument[topologies-]{topologies}
\externaldocument[descent-]{descent}
\externaldocument[perfect-]{perfect}
\externaldocument[more-morphisms-]{more-morphisms}
\externaldocument[flat-]{flat}
\externaldocument[groupoids-]{groupoids}
\externaldocument[more-groupoids-]{more-groupoids}
\externaldocument[etale-]{etale}
\externaldocument[chow-]{chow}
\externaldocument[intersection-]{intersection}
\externaldocument[pic-]{pic}
\externaldocument[weil-]{weil}
\externaldocument[adequate-]{adequate}
\externaldocument[dualizing-]{dualizing}
\externaldocument[duality-]{duality}
\externaldocument[discriminant-]{discriminant}
\externaldocument[derham-]{derham}
\externaldocument[local-cohomology-]{local-cohomology}
\externaldocument[algebraization-]{algebraization}
\externaldocument[curves-]{curves}
\externaldocument[resolve-]{resolve}
\externaldocument[models-]{models}
\externaldocument[functors-]{functors}
\externaldocument[equiv-]{equiv}
\externaldocument[pione-]{pione}
\externaldocument[etale-cohomology-]{etale-cohomology}
\externaldocument[proetale-]{proetale}
\externaldocument[relative-cycles-]{relative-cycles}
\externaldocument[more-etale-]{more-etale}
\externaldocument[trace-]{trace}
\externaldocument[crystalline-]{crystalline}
\externaldocument[spaces-]{spaces}
\externaldocument[spaces-properties-]{spaces-properties}
\externaldocument[spaces-morphisms-]{spaces-morphisms}
\externaldocument[decent-spaces-]{decent-spaces}
\externaldocument[spaces-cohomology-]{spaces-cohomology}
\externaldocument[spaces-limits-]{spaces-limits}
\externaldocument[spaces-divisors-]{spaces-divisors}
\externaldocument[spaces-over-fields-]{spaces-over-fields}
\externaldocument[spaces-topologies-]{spaces-topologies}
\externaldocument[spaces-descent-]{spaces-descent}
\externaldocument[spaces-perfect-]{spaces-perfect}
\externaldocument[spaces-more-morphisms-]{spaces-more-morphisms}
\externaldocument[spaces-flat-]{spaces-flat}
\externaldocument[spaces-groupoids-]{spaces-groupoids}
\externaldocument[spaces-more-groupoids-]{spaces-more-groupoids}
\externaldocument[bootstrap-]{bootstrap}
\externaldocument[spaces-pushouts-]{spaces-pushouts}
\externaldocument[spaces-chow-]{spaces-chow}
\externaldocument[groupoids-quotients-]{groupoids-quotients}
\externaldocument[spaces-more-cohomology-]{spaces-more-cohomology}
\externaldocument[spaces-simplicial-]{spaces-simplicial}
\externaldocument[spaces-duality-]{spaces-duality}
\externaldocument[formal-spaces-]{formal-spaces}
\externaldocument[restricted-]{restricted}
\externaldocument[spaces-resolve-]{spaces-resolve}
\externaldocument[formal-defos-]{formal-defos}
\externaldocument[defos-]{defos}
\externaldocument[cotangent-]{cotangent}
\externaldocument[examples-defos-]{examples-defos}
\externaldocument[algebraic-]{algebraic}
\externaldocument[examples-stacks-]{examples-stacks}
\externaldocument[stacks-sheaves-]{stacks-sheaves}
\externaldocument[criteria-]{criteria}
\externaldocument[artin-]{artin}
\externaldocument[quot-]{quot}
\externaldocument[stacks-properties-]{stacks-properties}
\externaldocument[stacks-morphisms-]{stacks-morphisms}
\externaldocument[stacks-limits-]{stacks-limits}
\externaldocument[stacks-cohomology-]{stacks-cohomology}
\externaldocument[stacks-perfect-]{stacks-perfect}
\externaldocument[stacks-introduction-]{stacks-introduction}
\externaldocument[stacks-more-morphisms-]{stacks-more-morphisms}
\externaldocument[stacks-geometry-]{stacks-geometry}
\externaldocument[moduli-]{moduli}
\externaldocument[moduli-curves-]{moduli-curves}
\externaldocument[examples-]{examples}
\externaldocument[exercises-]{exercises}
\externaldocument[guide-]{guide}
\externaldocument[desirables-]{desirables}
\externaldocument[coding-]{coding}
\externaldocument[obsolete-]{obsolete}
\externaldocument[fdl-]{fdl}
\externaldocument[index-]{index}

% Theorem environments.
%
\theoremstyle{plain}
\newtheorem{theorem}[subsection]{Theorem}
\newtheorem{proposition}[subsection]{Proposition}
\newtheorem{lemma}[subsection]{Lemma}

\theoremstyle{definition}
\newtheorem{definition}[subsection]{Definition}
\newtheorem{example}[subsection]{Example}
\newtheorem{exercise}[subsection]{Exercise}
\newtheorem{situation}[subsection]{Situation}

\theoremstyle{remark}
\newtheorem{remark}[subsection]{Remark}
\newtheorem{remarks}[subsection]{Remarks}

\numberwithin{equation}{subsection}

% Macros
%
\def\lim{\mathop{\mathrm{lim}}\nolimits}
\def\colim{\mathop{\mathrm{colim}}\nolimits}
\def\Spec{\mathop{\mathrm{Spec}}}
\def\Hom{\mathop{\mathrm{Hom}}\nolimits}
\def\Ext{\mathop{\mathrm{Ext}}\nolimits}
\def\SheafHom{\mathop{\mathcal{H}\!\mathit{om}}\nolimits}
\def\SheafExt{\mathop{\mathcal{E}\!\mathit{xt}}\nolimits}
\def\Sch{\mathit{Sch}}
\def\Mor{\mathop{\mathrm{Mor}}\nolimits}
\def\Ob{\mathop{\mathrm{Ob}}\nolimits}
\def\Sh{\mathop{\mathit{Sh}}\nolimits}
\def\NL{\mathop{N\!L}\nolimits}
\def\CH{\mathop{\mathrm{CH}}\nolimits}
\def\proetale{{pro\text{-}\acute{e}tale}}
\def\etale{{\acute{e}tale}}
\def\QCoh{\mathit{QCoh}}
\def\Ker{\mathop{\mathrm{Ker}}}
\def\Im{\mathop{\mathrm{Im}}}
\def\Coker{\mathop{\mathrm{Coker}}}
\def\Coim{\mathop{\mathrm{Coim}}}

% Boxtimes
%
\DeclareMathSymbol{\boxtimes}{\mathbin}{AMSa}{"02}

%
% Macros for moduli stacks/spaces
%
\def\QCohstack{\mathcal{QC}\!\mathit{oh}}
\def\Cohstack{\mathcal{C}\!\mathit{oh}}
\def\Spacesstack{\mathcal{S}\!\mathit{paces}}
\def\Quotfunctor{\mathrm{Quot}}
\def\Hilbfunctor{\mathrm{Hilb}}
\def\Curvesstack{\mathcal{C}\!\mathit{urves}}
\def\Polarizedstack{\mathcal{P}\!\mathit{olarized}}
\def\Complexesstack{\mathcal{C}\!\mathit{omplexes}}
% \Pic is the operator that assigns to X its picard group, usage \Pic(X)
% \Picardstack_{X/B} denotes the Picard stack of X over B
% \Picardfunctor_{X/B} denotes the Picard functor of X over B
\def\Pic{\mathop{\mathrm{Pic}}\nolimits}
\def\Picardstack{\mathcal{P}\!\mathit{ic}}
\def\Picardfunctor{\mathrm{Pic}}
\def\Deformationcategory{\mathcal{D}\!\mathit{ef}}
\usepackage{cleveref}
\crefname{originallemma}{original lemma}{original lemmas}
\Crefname{originallemma}{Original lemma}{Original lemmas}
\makeatletter
\newlabel{lemma-etale-etale-local-source-target}{{0CEZ}{source}{}{sourcepointer-1}{}}
\newlabel{lemma-etale-etale-local-source-target@cref}{{[originallemma][0][]0CEZ}{[1][0][]source}{}{}{}}
\makeatother
\crefname{originallemma}{original lemma}{original lemmas}
\Crefname{originallemma}{Original lemma}{Original lemmas}
\makeatletter
\newlabel{lemma-etale-local-source-target}{{04R2}{source}{}{sourcepointer-2}{}}
\newlabel{lemma-etale-local-source-target@cref}{{[originallemma][0][]04R2}{[1][0][]source}{}{}{}}
\makeatother
\crefname{originallemma}{original lemma}{original lemmas}
\Crefname{originallemma}{Original lemma}{Original lemmas}
\makeatletter
\newlabel{lemma-largest-open-of-the-source}{{0CEY}{source}{}{sourcepointer-3}{}}
\newlabel{lemma-largest-open-of-the-source@cref}{{[originallemma][0][]0CEY}{[1][0][]source}{}{}{}}
\makeatother
\crefname{originallemma}{original lemma}{original lemmas}
\Crefname{originallemma}{Original lemma}{Original lemmas}
\makeatletter
\newlabel{lemma-orbits}{{0CF0}{source}{}{sourcepointer-4}{}}
\newlabel{lemma-orbits@cref}{{[originallemma][0][]0CF0}{[1][0][]source}{}{}{}}
\makeatother
\crefname{originallemma}{original lemma}{original lemmas}
\Crefname{originallemma}{Original lemma}{Original lemmas}
\makeatletter
\newlabel{morphisms-lemma-smooth-etale-over-affine-space}{{054L}{source}{}{sourcepointer-5}{}}
\newlabel{morphisms-lemma-smooth-etale-over-affine-space@cref}{{[originallemma][0][]054L}{[1][0][]source}{}{}{}}
\makeatother
\title{Etale locality on the source plus smooth locality on the target implies locality on the source-and-target}
\author{The Stacks Project Authors}
\date{stacks-project:descent.tex:local-snapshot-20261012; GNU FDL 1.2 or later}
\begin{document}
\maketitle
\noindent\textit{Editorial note.} All delivered mathematical text is the author's own, copied from the pinned Stacks Project chapter file descent.tex; no statement or proof step was rewritten, reordered or removed.\par\medskip
\noindent\textit{Editorial note.} The proof contains the author's own remark that the naive translation-invariance argument would require extra etale localization on the target and would become annoying, followed by the choice to use the automorphisms of the group action lemma instead; the remark is reproduced verbatim and is part of the argument's structure.\par\medskip
\noindent\textit{Editorial note.} The proof also contains the parenthetical hint '(Details omitted; hint: compute the derivatives.)' for the etaleness of the p-th power map; the parenthetical is reproduced verbatim and the surrounding argument is complete as written.\par\medskip
\noindent\textit{Editorial note.} The counted spans carry 5 navigation stubs; one of them points at the published item of this chapter whose label is lemma-orbits, which is used here as a prerequisite rather than restated, so this card is not a duplicate of it.\par\medskip
\noindent\textit{Editorial note.} The preamble span points at the chapter's own preamble.tex shipped with the frozen source, used verbatim as staged.\par\medskip
\noindent\textit{Allowed standard prerequisites (editorial).}\begin{itemize}
\item Morphisms of germs attached to a property and the locus on which the property holds. Same chapter, Lemma lemma-largest-open-of-the-source, named in the statement at line 7931. Supplies the object W(-) being compared.
\item Smooth morphisms are etale locally standard smooth over an affine space. Morphisms, Lemma morphisms-lemma-smooth-etale-over-affine-space, invoked at proof lines 7950-7952. Supplies the etale morphism of the target to an affine space over the base.
\item Etale locality on the source implies etale locality on the source-and-target, and the comparison of loci along the resulting square. Same chapter, Lemmas lemma-etale-local-source-target and lemma-etale-etale-local-source-target, invoked at proof lines 7961-7966 and 8033. Supplies the identification of the loci and the invariance statement used with the Frobenius-like map.
\item The group of elementary automorphisms of affine space and the structure of its invariant open subschemes. Same chapter, Lemma lemma-orbits, invoked at proof lines 7981 and 8006-8011. Supplies the description of the invariant open in each fibre, which is then excluded in characteristic p.
\end{itemize}\medskip
\ifdefined\hypertarget\hypertarget{sourcepointer-1}{}\fi
\noindent\textit{Source reference.} Original lemma 0CEZ: sources/stacks-descent/pinned-original-source.tex, source lines 7787--7787. Pointer outside the counted span, invoked once for context; the load-bearing argument of the counted statement is the author's proof reproduced in full in the proof span.\par\medskip
\ifdefined\hypertarget\hypertarget{sourcepointer-2}{}\fi
\noindent\textit{Source reference.} Original lemma 04R2: sources/stacks-descent/pinned-original-source.tex, source lines 7532--7532. Pointer outside the counted span, invoked once for context; the load-bearing argument of the counted statement is the author's proof reproduced in full in the proof span.\par\medskip
\ifdefined\hypertarget\hypertarget{sourcepointer-3}{}\fi
\noindent\textit{Source reference.} Original lemma 0CEY: sources/stacks-descent/pinned-original-source.tex, source lines 6890--6890. Pointer outside the counted span, invoked once for context; the load-bearing argument of the counted statement is the author's proof reproduced in full in the proof span.\par\medskip
\ifdefined\hypertarget\hypertarget{sourcepointer-4}{}\fi
\noindent\textit{Source reference.} Original lemma 0CF0: sources/stacks-descent/pinned-original-source.tex, source lines 7847--7847. Pointer outside the counted span, invoked once for context; the load-bearing argument of the counted statement is the author's proof reproduced in full in the proof span.\par\medskip
\ifdefined\hypertarget\hypertarget{sourcepointer-5}{}\fi
\noindent\textit{Source reference.} Original lemma 054L: sources/stacks-chow/tags.txt, source lines 6644--6644. Pointer outside the counted span, invoked once for context; the load-bearing argument of the counted statement is the author's proof reproduced in full in the proof span.\par\medskip
\begin{lemma}
\label{lemma-etale-tau-local-source-target}
Let $\mathcal{P}$ be a property of morphisms of schemes. Assume
\begin{enumerate}
\item $\mathcal{P}$ is \'etale local on the source,
\item $\mathcal{P}$ is smooth local on the target,
\item $\mathcal{P}$ is stable under postcomposing with open immersions:
if $f : X \to Y$ has $\mathcal{P}$ and $Y \subset Z$ is an open
subscheme then $X \to Z$ has $\mathcal{P}$.
\end{enumerate}
Given a commutative diagram of schemes
$$
\vcenter{
\xymatrix{
X' \ar[d]_{g'} \ar[r]_{f'} & Y' \ar[d]^g \\
X \ar[r]^f & Y
}
}
\quad\text{with points}\quad
\vcenter{
\xymatrix{
x' \ar[d] \ar[r] & y' \ar[d] \\
x \ar[r] & y
}
}
$$
such that $g$ is smooth $y'$ and $X' \to X \times_Y Y'$ is \'etale
at $x'$, then $x \in W(f) \Leftrightarrow x' \in W(f')$
where $W(-)$ is as in Lemma \ref{lemma-largest-open-of-the-source}.
\end{lemma}

\begin{proof}
Since $\mathcal{P}$ is \'etale local on the source we see
that $x \in W(f)$ if and only if the image of $x$ in
$X \times_Y Y'$ is in $W(X \times_Y Y' \to Y')$. Hence we
may assume the diagram in the lemma is cartesian.

\medskip\noindent
Assume $x \in W(f)$. Since $\mathcal{P}$ is smooth local on the target
we see that $(g')^{-1}W(f) = W(f) \times_Y Y' \to Y'$ has $\mathcal{P}$.
Hence $(g')^{-1}W(f) \subset W(f')$. We conclude $x' \in W(f')$.

\medskip\noindent
Assume $x' \in W(f')$.
For any open neighbourhood $V' \subset Y'$ of $y'$ we may replace
$Y'$ by $V'$ and $X'$ by $U' = (f')^{-1}V'$ because $V' \to Y'$ is smooth
and hence the base change $W(f') \cap U' \to V'$ of $W(f') \to Y'$
has property $\mathcal{P}$. Thus we may assume there exists
an \'etale morphism $Y' \to \mathbf{A}^n_Y$ over $Y$, see
Morphisms, Lemma \ref{morphisms-lemma-smooth-etale-over-affine-space}.
Picture
$$
\xymatrix{
X' \ar[r] \ar[d] & Y' \ar[d] \\
\mathbf{A}^n_X \ar[r]_{f_n} \ar[d] & \mathbf{A}^n_Y \ar[d] \\
X \ar[r]^f & Y
}
$$
By Lemma \ref{lemma-etale-local-source-target}
(and because \'etale coverings are smooth coverings)
we see that $\mathcal{P}$ is \'etale local on the source-and-target.
By Lemma \ref{lemma-etale-etale-local-source-target}
we see that $W(f')$ is the inverse image of
the open $W(f_n) \subset \mathbf{A}^n_X$. In particular
$W(f_n)$ contains a point lying over $x$.
After replacing $X$ by the image of $W(f_n)$ (which is open)
we may assume $W(f_n) \to X$ is surjective.
Claim: $W(f_n) = \mathbf{A}^n_X$.
The claim implies $f$ has $\mathcal{P}$ as
$\mathcal{P}$ is local in the smooth topology
and $\{\mathbf{A}^n_Y \to Y\}$ is a smooth covering.

\medskip\noindent
Essentially, the claim follows as $W(f_n) \subset \mathbf{A}^n_X$ is a
``translation invariant'' open which meets every fibre of
$\mathbf{A}^n_X \to X$. However, to produce an argument along these lines
one has to do \'etale localization on $Y$ to produce enough translations
and it becomes a bit annoying. Instead we use the automorphisms
of Lemma \ref{lemma-orbits} and \'etale morphisms of affine spaces.
We may assume $n \geq 2$. Namely, if $n = 0$, then we are done.
If $n = 1$, then we consider the diagram
$$
\xymatrix{
\mathbf{A}^2_X \ar[r]_{f_2} \ar[d]_p & \mathbf{A}^2_Y \ar[d] \\
\mathbf{A}^1_X \ar[r]^{f_1} & \mathbf{A}^1_Y
}
$$
We have $p^{-1}(W(f_1)) \subset W(f_2)$ (see first paragraph
of the proof). Thus $W(f_2) \to X$ is still surjective
and we may work with $f_2$. Assume $n \geq 2$.

\medskip\noindent
For any $1 \leq i, j \leq n$ with $i \not = j$ and $d \geq 0$
denote $T_{i, j, d}$ the automorphism of $\mathbf{A}^n$ defined
in Lemma \ref{lemma-orbits}. Then we get a commutative diagram
$$
\xymatrix{
\mathbf{A}^n_X \ar[r]_{f_n} \ar[d]_{T_{i, j, d}} &
\mathbf{A}^n_Y \ar[d]^{T_{i, j, d}} \\
\mathbf{A}^n_X \ar[r]^{f_n} & \mathbf{A}^n_Y
}
$$
whose vertical arrows are isomorphisms. We conclude that
$T_{i, j, d}(W(f_n)) = W(f_n)$. Applying Lemma \ref{lemma-orbits}
we conclude for any $x \in X$ the fibre $W(f_n)_x \subset \mathbf{A}^n_x$ is
either $\mathbf{A}^n_x$ (this is what we want) or $\kappa(x)$
has characteristic $p > 0$ and $W(f_n)_x$
is the complement of a finite set $Z_x \subset \mathbf{A}^n_x$
of closed points. The second possibility cannot occur. Namely,
consider the morphism $T_p : \mathbf{A}^n \to \mathbf{A}^n$ given by
$$
(x_1, \ldots, x_n) \mapsto (x_1 - x_1^p, \ldots, x_n - x_n^p)
$$
As above we get a commutative diagram
$$
\xymatrix{
\mathbf{A}^n_X \ar[r]_{f_n} \ar[d]_{T_p} &
\mathbf{A}^n_Y \ar[d]^{T_p} \\
\mathbf{A}^n_X \ar[r]^{f_n} & \mathbf{A}^n_Y
}
$$
The morphism $T_p : \mathbf{A}^n_X \to \mathbf{A}^n_X$
is \'etale at every point lying over $x$
and the morphism $T_p : \mathbf{A}^n_Y \to \mathbf{A}^n_Y$
is \'etale at every point lying over the image of $x$ in $Y$.
(Details omitted; hint: compute the derivatives.)
We conclude that
$$
T_p^{-1}(W) \cap \mathbf{A}^n_x = W \cap \mathbf{A}^n_x
$$
by Lemma \ref{lemma-etale-etale-local-source-target}
(we've already seen $\mathcal{P}$ is
\'etale local on the source-and-target).
Since $T_p : \mathbf{A}^n_x \to \mathbf{A}^n_x$ is finite \'etale
of degree $p^n > 1$ we see that if $Z_x$ is not empty then it contains
$T_p^{-1}(Z_x)$ which is bigger. This contradiction finishes
the proof.
\end{proof}

\end{document}
